\section{Introduction}
\label{sec:introduction}

Heuristic evaluation is one of the most widely used inspection methods in human--computer interaction \cite{nielsen1990heuristic,nielsen1994heuristic}. Teaching it begins with helping learners recognise each heuristic in real interfaces, and this is where novices most often struggle: they find fewer problems than specialists and report difficulty applying abstract definitions to concrete cases \cite{nielsen1992finding,abulfaraj2020detailed}. Reading each heuristic in isolation rarely gives enough anchoring for a learner to identify the same idea when it appears in a screen. The early teaching of Nielsen's heuristics is therefore a practical challenge for HCI educators.

The Brazilian HCI community has responded to this kind of challenge with game-based and playful pedagogy \cite{silveira2024praticas}. Experience reports describe gamified courses \cite{miranda2019gamificacao}, card games \cite{darin2019goople}, serious games \cite{sales2020jogosserios,guimaraes2022avaliacao}, and other playful artefacts \cite{santos2025ihcweek,rodrigues2025dossies} used to teach different HCI topics. For Nielsen's heuristics specifically, the closest published artefact in the community is a physical board game \cite{geremias2022desvendando}; web-based, scenario-driven alternatives are still rare.

This paper reports the design and classroom evaluation of a web-based educational game that supports the understanding of Nielsen's ten heuristics through paired scenarios.\footnote{Link omitted for double-blind review.} For each heuristic, the player completes the same concrete task in two interfaces: one that violates the heuristic and one that respects it. A narrator frames each scenario before and after the interaction. Light gamification --- stars, a progress sidebar, and a celebratory closing screen that greets the player as a ``heuristics expert'' after the tenth scenario --- helps the player work through the full sequence in a single session.

The game was applied in an HCI evaluation course at a Brazilian university (omitted for double-blind review), and further participants were reached through the research team's network. After playing, participants answered a questionnaire with five categorical demographic items and five open-ended questions about their experience. The open-ended responses were analysed through qualitative content analysis \cite{bardin2011}, under an ethics protocol approved by the local research ethics committee (CAAE omitted for double-blind review).

This experience report contributes (a) the design and implementation of the artefact, (b) a replicable classroom procedure for using it in an HCI course, and (c) a qualitative reading of how participants described the support the game provided for understanding Nielsen's heuristics. The remainder of the paper is organised as follows. Section~\ref{sec:background} reviews the theoretical background. Section~\ref{sec:related} discusses related work. Section~\ref{sec:game} describes the game. Section~\ref{sec:method} presents the methodology and Section~\ref{sec:ethics} the ethical care taken. Sections~\ref{sec:results} and~\ref{sec:discussion} present and discuss the results. Section~\ref{sec:conclusion} closes with final considerations.
